% ============================================================================
% Chapter 2 --- Conceptual frame
% ----------------------------------------------------------------------------
% Purpose : establish a shared mental model of the TAC (Trusted Activity
%           Context) and the perimeter of the document. Provides:
%             §2.1 formal definition of TAC for the standalone reader;
%             §2.2 end-to-end activation sequence with figure placeholder;
%             §2.3 endpoint OS integration matrix;
%             §2.4 lifecycle state machine of a TAC;
%             §2.5 out-of-scope statement.
% Page target: ~10 pages (revised upward from ~6 after F3 decision).
% Dependencies: introduces concepts used by every subsequent chapter.
% Cross-document: the formal TAC definition aligns with Section 3.6 of the
%                 main report; the technical reference's standalone status
%                 means the definition is fully reproduced here, not recapped.
% ============================================================================

\chapter{Conceptual frame}
\label{ch:conceptual-frame}

This chapter establishes the conceptual material that the rest of
the Technical Reference applies. It develops the formal definition
of the Trusted Activity Context (TAC), the end-to-end activation
sequence that turns a TAC declaration into running infrastructure,
the endpoint operating-system integration matrix, and the
lifecycle state machine that governs a TAC from draft to archival.
A reader who has consulted the companion main report will find
parts of this material familiar; the chapter is nevertheless
written to be self-contained for the standalone reader.

%-----------------------------------------------------------------------------
\section{The Trusted Activity Context concept}
\label{sec:cf-tac-concept}

\paragraph{Definition.} A \keyword{Trusted Activity Context}
(\gls{tac}) is a coherent, validated environment that assembles
the data, tools, services, and access rights required for a
specific activity --- a research project, a digital examination,
an administrative process, a teaching session --- under
documented institutional governance. Within a TAC, the
identity of subjects, the classification of resources, the
network attachment, the policy decisions, the cryptographic
posture, and the audit trail are aligned with the activity's
objectives and risk level, and are operationally bounded by the
TAC's declared scope.

\paragraph{Why ``trusted''.} The qualifier follows the
dependability tradition of Avizienis et al.~\cite{avizienis2004}:
an entity is \emph{trusted} when it delivers service that can
justifiably be relied upon. The trust in a TAC is consequently not
assumed; it is the result of the contractual obligations placed on
each layer by the architecture of Part~II and verified by the
conformance suite of \trchap{ch:conformance}{}. A TAC is trusted
because its properties are checkable.

\paragraph{Why ``activity''.} The qualifier is the central
innovation. Existing access-control and isolation frameworks are
typically organised around either subjects (who a user is) or
resources (what a service contains). A TAC is organised around
\emph{what is being done}: the same subject, holding the same
attributes, accessing the same resource, may operate in a
different TAC depending on the activity engaged. This re-frames
several recurring difficulties in academic IT --- the
researcher who is also a teacher, the visiting collaborator who
is also a host, the administrative process that touches research
data --- as switches between TACs rather than as exceptions to a
single role-based model.

\paragraph{Why ``context''.} A TAC is not a single object; it is
the assembled configuration that an activity requires. It
includes:

\begin{itemize}[leftmargin=1.5em,nosep]
  \item a \emph{subject set}, typically with specific role
    activations rather than blanket role memberships;
  \item a \emph{resource set} --- data, tools, services ---
    consistent with the activity's classification;
  \item a \emph{network attachment} appropriate to the activity's
    zone (research, teaching, administration, sensitive);
  \item a \emph{policy envelope} that admits the operations the
    activity requires and denies those it does not;
  \item a \emph{cryptographic posture} (key custody, certificate
    chain) consistent with the resource classification;
  \item a \emph{logging and audit trail} that allows every
    consequential decision under the TAC to be reconstructed.
\end{itemize}

The TAC is the unit of \emph{governance} that brings these
elements together, so that they can be specified, reviewed,
provisioned, audited, and decommissioned as a coherent whole
rather than as a coincidence of separately managed
configurations.

\paragraph{Disambiguation.} The term \emph{Trusted Activity
Context} should not be confused with IBM~DB2's identically named
database connection feature, nor with the Trusted Computing
concepts of the Trusted Computing Group. Here the term has the
specific meaning developed above.

\paragraph{Position relative to existing frameworks.} The TAC
concept draws on and extends several established frameworks
without being identical to any of them. Role-based access
control (RBAC)~\cite{sandhu-rbac-1996} provides session-bounded
role activation but does not assemble resources, network, and
audit into a single governance unit. Task-Based Authorization
Controls (TBAC)~\cite{thomas-sandhu-tbac-1997} tie permissions to
task lifecycles and supply part of the lifecycle intuition but do
not extend to the cross-layer contractual obligations.
Attribute-Based Access Control (ABAC, NIST SP
800-162)~\cite{nist-sp800-162} provides the per-decision evaluation
model that the TAC policy envelope uses, but does not extend to
the governance-level abstraction. The contextual-awareness
technique of NIST SP 800-160 Vol.~2 names contextual factors as
inputs to security decisions but does not propose a unit of
operational governance over those inputs. Workspace and tenancy
products (cloud workspaces, Kubernetes namespaces with RBAC)
operationalise aspects of the pattern but at the level of a
single platform, not at the level of an institutional contract
spanning identity, network, compute, storage, observability and
governance.

\paragraph{Four founding principles.} A TAC, in this Reference's
sense, satisfies four properties:

\begin{description}[leftmargin=2em,nosep,style=nextline,font=\bfseries]
  \item[Trustworthiness] Each context is technically bounded
    and monitored so that the institution can grant a high
    degree of autonomy within it without surrendering control
    of the institutional perimeter.
  \item[Dependability] Failures, misconfigurations, or security
    incidents within one context remain contained and do not
    compromise the institutional core or other contexts.
  \item[User-centricity] The digital environment adapts to the
    user's current activity rather than imposing a uniform
    security model on every situation. Restriction is the
    consequence of activity-specific risk, not of
    institutional default.
  \item[Observability] Every action within a TAC is logged and
    reconstructable, so that autonomy is accompanied by
    accountability.
\end{description}

\paragraph{Isolation as enabling, not restrictive.} The TAC
model reverses the conventional framing of isolation. In the
conventional framing, isolation \emph{restricts} users; in the
TAC model, isolation \emph{enables} autonomy by guaranteeing
that mistakes, experiments, and attacks within a context cannot
escape it. The forestry analogy applies: a forest is not made
safer by planting trees in a single continuous mass; the
\emph{construction of firebreaks} is what prevents a blaze from
spreading to the entire woodland. TACs are the institutional
firebreaks of the digital infrastructure.

\paragraph{Relation to segmentation.} A mature institution
combines a relatively stable network segmentation (zones, as
described in \trchap{ch:network}{}) with multiple dynamic TACs
that operate within those zones. Segmentation is spatial and
relatively stable; TACs are activity-bound and relatively
dynamic. The two mechanisms are complementary: segmentation
limits where flows can circulate; TACs determine in what
validated environment an activity proceeds.

%-----------------------------------------------------------------------------
\section{End-to-end activation sequence}
\label{sec:cf-activation-sequence}

A TAC is realised by the cross-layer cooperation of the
architecture of Part~II. The activation sequence below describes
what happens, end to end, when a subject activates a TAC for a
specific activity. The sequence is presented at the level of
named functional roles; the actual products that play these
roles depend on the institution's chosen instantiation
(\trchap{ch:identity}{} through \trchap{ch:orchestration}{}).

\paragraph{Functional roles in the sequence.} The sequence
involves: \textsc{Identity} (the institution's identity provider,
\trchap{ch:identity}{}); \textsc{Policy} (the policy decision
point, \trchap{ch:policy}{}); \textsc{Network} (the network
attachment and segmentation infrastructure,
\trchap{ch:network}{}); \textsc{Compute}, \textsc{Storage}, and
\textsc{Endpoint} (the execution-plane layers,
\trchap{ch:compute}{}--\trchap{ch:endpoint}{});
\textsc{Crypto} (cryptographic services,
\trchap{ch:crypto}{}); \textsc{Audit} (the observability layer,
\trchap{ch:observability}{}); \textsc{Orchestration} (the
cross-layer glue, \trchap{ch:orchestration}{}).

\paragraph{The activation sequence.} The numbered steps below
correspond to a typical TAC activation. Steps may overlap in time
in a real implementation; the numbering reflects logical order,
not strict serial execution.

\begin{enumerate}[leftmargin=2em,nosep,label=\arabic*.]
  \item \textsc{Subject authentication.} The subject authenticates
    against \textsc{Identity}; a short-lived assertion is issued
    carrying the institutional attribute schema.
  \item \textsc{Activity declaration.} The subject (directly, or
    through an institutional service such as a research-workspace
    portal) requests activation of a named TAC for a stated
    activity.
  \item \textsc{Policy evaluation.} \textsc{Policy} receives the
    activation request, combines the subject attributes with the
    TAC declaration (resource set, classification, environment),
    and produces a decision. The decision is one of \texttt{allow},
    \texttt{deny}, or \texttt{indeterminate}, accompanied by
    structured obligations (session timeout, audit class,
    additional factor required).
  \item \textsc{Orchestration triggers.} On \texttt{allow},
    \textsc{Orchestration} receives the decision and dispatches
    cross-layer commands to realise the TAC.
  \item \textsc{Network attachment.} \textsc{Network} attaches
    the subject's session to the appropriate zone, with
    per-session attributes (zone, egress allowance, rate limits)
    derived from the policy decision.
  \item \textsc{Compute provisioning.} \textsc{Compute}
    provisions or attaches the workload context (a virtual
    machine, a container namespace, a research workspace) with
    the institutional schema's attributes propagated.
  \item \textsc{Storage binding.} \textsc{Storage} binds the
    appropriate volumes, applying classification-derived access
    controls and ensuring keys are obtained from
    \textsc{Crypto} consistent with the resource classification.
  \item \textsc{Endpoint posture verification.} \textsc{Endpoint}
    confirms that the subject's device meets the posture
    requirements declared by the policy decision; mismatch is
    fed back to \textsc{Policy} for re-evaluation or
    quarantine.
  \item \textsc{Audit correlation.} \textsc{Audit} receives
    events from every layer with the same TAC correlation
    identifier, allowing the activation chain to be reconstructed
    end-to-end.
  \item \textsc{Activity proceeds.} The subject conducts the
    activity within the bounded context. Each subsequent
    consequential operation under the TAC produces events
    correlated to the same identifier.
  \item \textsc{Deactivation.} On the activity's natural end
    (timeout, explicit deactivation, lifecycle transition;
    \S\ref{sec:cf-lifecycle}), \textsc{Orchestration} reverses
    the consequential changes. \textsc{Audit} records the
    deactivation; \textsc{Storage} preserves classification
    metadata on persistent artefacts; \textsc{Network} releases
    the attachment; \textsc{Compute} releases the workload
    context.
\end{enumerate}

\begin{figure}[H]
\centering
\begin{tikzpicture}[
    every node/.style = {font=\footnotesize},
    role/.style = {rectangle, rounded corners=2pt,
                   draw=NavyBlue!70!black, fill=NavyBlue!10,
                   minimum width=15mm, minimum height=7mm,
                   line width=0.4pt, align=center, font=\footnotesize\bfseries},
    msg/.style = {-{Stealth[length=1.8mm]}, draw=NavyBlue!60!black,
                  line width=0.4pt},
    return/.style = {-{Stealth[length=1.8mm]}, draw=Slate!60!black,
                     line width=0.35pt, dashed},
    lifeline/.style = {draw=Slate!50, line width=0.3pt, dashed}
]
% Role headers (top)
\foreach \name/\x in {Subject/0, Identity/1.5, Policy/3.2,
                       Orchestr./4.9, Network/6.6, Compute/8.3,
                       Storage/9.8, Endpoint/11.5, Audit/13.2} {
  \node[role] (\name) at (\x,0) {\name};
}

% Lifelines
\foreach \x in {0,1.5,3.2,4.9,6.6,8.3,9.8,11.5,13.2} {
  \draw[lifeline] (\x,-0.4) -- (\x,-14.5);
}

% Messages (numbered 1-11) -- 50% increased vertical spacing
\draw[msg] (0,-1.0) -- node[above,font=\scriptsize]{1.\,authenticate} (1.5,-1.0);
\draw[msg] (0,-1.9) -- node[above,font=\scriptsize]{2.\,activate TAC} (3.2,-1.9);
\draw[msg] (3.2,-2.8) -- node[above,font=\scriptsize]{3.\,decision} node[below,font=\scriptsize]{(allow + obligations)} (4.9,-2.8);
\draw[msg] (4.9,-4.0) -- node[above,font=\scriptsize]{4.\,attach} (6.6,-4.0);
\draw[msg] (4.9,-4.9) -- node[above,font=\scriptsize]{5.\,provision} (8.3,-4.9);
\draw[msg] (4.9,-5.8) -- node[above,font=\scriptsize]{6.\,bind volumes} (9.8,-5.8);
\draw[msg] (4.9,-6.7) -- node[above,font=\scriptsize]{7.\,verify posture} (11.5,-6.7);
\draw[return] (11.5,-7.6) -- node[above,font=\scriptsize]{posture OK} (3.2,-7.6);
\draw[msg] (6.6,-8.5) to[bend right=10] node[below,font=\scriptsize]{8.\,audit events} (13.2,-8.5);
\draw[msg] (8.3,-9.4) to[bend right=8]  (13.2,-9.4);
\draw[msg] (9.8,-10.3) to[bend right=8]  (13.2,-10.3);
\draw[msg] (11.5,-11.2) -- (13.2,-11.2);

% Activity proceeds (loop)
\draw[draw=Slate!60, line width=0.3pt]
  (0,-11.8) rectangle (13.5,-12.7);
\node[font=\scriptsize\itshape, anchor=west] at (0.1,-12.25) {%
  9.--10.\,Activity proceeds; subsequent operations correlate to same
  \texttt{correlation\_id}; events emitted to Audit};

% Deactivation
\draw[msg] (0,-13.3) -- node[above,font=\scriptsize]{11.\,deactivate} (4.9,-13.3);
\draw[return] (4.9,-14.2) -- node[above,font=\scriptsize\itshape,
              color=Slate!70!black]{reverse: release attachment, workload, volumes; archive}
              (13.2,-14.2);

% Correlation_id thread (vertical Crimson line)
\draw[Crimson!70!black, line width=0.8pt, dotted]
  (-0.5,-1.0) -- (-0.5,-14.2);
\node[Crimson!70!black, font=\scriptsize\bfseries, rotate=90, anchor=south]
  at (-0.5,-7.5) {TAC correlation\_id};
\end{tikzpicture}
\caption{End-to-end TAC activation sequence: eleven steps across
the subject and the principal functional-role lifelines
(cryptographic services participate within the certificate- and
key-acquisition steps rather than as a separate lifeline). The TAC \texttt{correlation\_id}
(Crimson) propagates through every step; deactivation reverses
every consequential change.}
\label{fig:cf-activation-sequence}
\end{figure}

\paragraph{What the sequence guarantees.} If every layer respects
the contracts of Part~II, the sequence produces three
guarantees: every TAC activation is the consequence of a recorded
policy decision; every consequential change is reversible by
the deactivation path; every layer's actions can be correlated
through the TAC identifier for forensic and audit reconstruction.
These guarantees are the operational expression of the four
founding principles of \S\ref{sec:cf-tac-concept}.

%-----------------------------------------------------------------------------
\section{Endpoint operating-system integration matrix}
\label{sec:cf-endpoint-matrix}

The TAC model relies on the endpoint to enforce a subset of
the policy decisions and to report posture to the policy plane.
The endpoint operating systems encountered at \loc{institution-type}
have differing capabilities and
differing maturity for TAC integration. The matrix in
Table~\ref{tab:cf-endpoint-matrix} sketches the principal
capabilities; \trchap{ch:endpoint}{} develops the operational
detail and the OSS-vs-commercial trade-offs.

\begin{table}[H]
\centering\small
\caption{Sketch of endpoint integration capabilities by operating
system. Each cell is indicative; specific institutional
deployments vary.}
\label{tab:cf-endpoint-matrix}
\begin{tabularx}{\textwidth}{%
  >{\hsize=0.9\hsize\linewidth=\hsize\bfseries\raggedright\arraybackslash}X%
  >{\hsize=1.0\hsize\linewidth=\hsize\raggedright\arraybackslash}X%
  >{\hsize=1.0\hsize\linewidth=\hsize\raggedright\arraybackslash}X%
  >{\hsize=1.1\hsize\linewidth=\hsize\raggedright\arraybackslash}X%
}
\toprule
Capability &
  \textbf{Linux (managed)} &
  \textbf{macOS (managed/BYOD)} &
  \textbf{Windows (corporate)} \\
\midrule
Federated identity &
  SSSD or realmd against AD/LDAP; OIDC clients for web apps &
  Apple Platform SSO with the IdP extension; Jamf~Connect as
  alternative &
  Entra~ID join with Conditional Access; Hello~for~Business \\
\addlinespace
Posture verification &
  Custom checks via \texttt{osquery}; package-manager state &
  MDM compliance profile (Jamf, Mosyle, Kandji); built-in
  managed-client signals &
  Intune compliance policies; Defender for Endpoint signals \\
\addlinespace
TAC-scoped tunnel &
  WireGuard or OpenVPN profile pushed by configuration management
  &
  Per-app VPN configuration profile delivered via MDM &
  Always-On VPN profile via Intune; SSE alternatives \\
\addlinespace
Configuration policy &
  Ansible, Puppet, Salt; declarative GitOps repositories &
  Munki, configuration profiles, Jamf, Apple Business Manager
  &
  GPO for AD environments; Intune for cloud-managed estates \\
\addlinespace
Audit ingestion &
  \texttt{auditd}, journald, syslog forwarder &
  Unified log forwarder (e.g.\ \texttt{log show} stream, EDR
  agent) &
  Windows Event Forwarding; EDR agent telemetry \\
\bottomrule
\end{tabularx}
\end{table}

\paragraph{Mobile platforms.} \institution{iOS} and
\institution{Android} share much with the
\institution{macOS}/\institution{Windows} columns through the
same MDM platforms (Jamf, Intune, Workspace ONE), with the
additional constraint that posture and configuration depend on
what each operating system exposes to enterprise management. As
\trchap{ch:endpoint}{} discusses explicitly, the open-source
ecosystem for iOS and Android at institutional scale is, at the
time of writing, less mature than for the desktop platforms; the
matrix above accordingly emphasises desktop coverage.

\paragraph{Other platforms.} \institution{BSD} and
\institution{illumos} environments, where present in the
research-computing footprint, can be brought into the matrix
through the same patterns as Linux for posture
(\texttt{osquery} has reasonable cross-platform coverage) and
through declarative configuration management. They are not
listed in the matrix because their institutional footprint is
typically small.

%-----------------------------------------------------------------------------
\section{Lifecycle of a TAC}
\label{sec:cf-lifecycle}

A TAC is not provisioned eternally; it has a lifecycle that the
institution governs. The state machine in
Figure~\ref{fig:cf-lifecycle} (and Table~\ref{tab:cf-lifecycle})
describes the principal states and the transitions between them.

\begin{table}[H]
\centering\small
\caption{TAC lifecycle: states and transition triggers.}
\label{tab:cf-lifecycle}
\begin{tabularx}{\textwidth}{%
  >{\hsize=0.5\hsize\linewidth=\hsize\bfseries\raggedright\arraybackslash}X%
  >{\hsize=1.5\hsize\linewidth=\hsize\raggedright\arraybackslash}X%
  >{\hsize=1.0\hsize\linewidth=\hsize\raggedright\arraybackslash}X%
}
\toprule
State & Meaning & Principal transitions out \\
\midrule
\textbf{Drafted}
  & The TAC declaration exists in the institutional registry but
    has not yet been validated. No infrastructure is provisioned.
  & to \emph{Validated} on review approval; to \emph{Discarded}
    on rejection. \\
\addlinespace
\textbf{Validated}
  & The TAC has passed institutional review (data classification
    confirmed, policy envelope reviewed, governance approval
    recorded). It can be activated on request.
  & to \emph{Provisioned} on first activation; to
    \emph{Decommissioned} if validity expires. \\
\addlinespace
\textbf{Provisioned}
  & The infrastructure to realise the TAC is in place but the
    activity is not currently running.
  & to \emph{Active} on activity start; to \emph{Suspended} on
    institutional pause. \\
\addlinespace
\textbf{Active}
  & The activity is in progress; subjects operate under the
    TAC; events are being correlated.
  & to \emph{Provisioned} (or \emph{Suspended}) on activity end. \\
\addlinespace
\textbf{Suspended}
  & The TAC is operationally paused: pending compliance review,
    stakeholder absence, or institutional pause. No activity
    proceeds.
  & to \emph{Active} on resumption; to \emph{Decommissioned} on
    permanent pause decision. \\
\addlinespace
\textbf{Decommissioned}
  & The infrastructure has been released; classification metadata
    on persistent artefacts is preserved per the institutional
    retention schedule.
  & to \emph{Archived} after the retention period expires. \\
\addlinespace
\textbf{Archived}
  & The TAC is closed. The audit trail and metadata remain
    available to institutional inquiry per retention policy; no
    further activity is possible.
  & terminal state. \\
\bottomrule
\end{tabularx}
\end{table}

\begin{figure}[H]
\centering
\begin{tikzpicture}[
    node distance = 14mm and 18mm,
    every node/.style = {font=\small},
    state/.style = {rectangle, rounded corners=2pt, draw=NavyBlue!70!black,
                    fill=NavyBlue!5, minimum width=24mm, minimum height=8mm,
                    align=center, line width=0.5pt},
    terminal/.style = {state, fill=Slate!10, draw=Slate!70!black},
    arrow/.style = {-{Stealth[length=2mm]}, draw=NavyBlue!60!black,
                    line width=0.4pt}
]
\node[state] (drafted)        {Drafted};
\node[state, right=of drafted] (validated)      {Validated};
\node[state, right=of validated] (provisioned)  {Provisioned};
\node[state, right=of provisioned] (active)     {Active};
\node[state, below=of provisioned] (suspended)  {Suspended};
\node[state, below=of validated] (decommissioned) {Decommissioned};
\node[terminal, below=of drafted] (archived)    {Archived};

\draw[arrow] (drafted)      -- node[above,font=\scriptsize]{review} (validated);
\draw[arrow] (validated)    -- node[above,font=\scriptsize]{first activation} (provisioned);
\draw[arrow] (provisioned)  -- node[above,font=\scriptsize]{activity start} (active);
\draw[arrow] (active.south west) to[bend left=10] node[below,font=\scriptsize]{activity end} (provisioned.south east);
\draw[arrow] (provisioned) -- node[right,font=\scriptsize]{pause} (suspended);
\draw[arrow] (active.south) to[bend left=20] (suspended.east);
\draw[arrow] (suspended.north west) to[bend left=15] node[left,font=\scriptsize]{resume} (provisioned.south west);
\draw[arrow] (suspended) -- node[above,font=\scriptsize]{permanent} (decommissioned);
\draw[arrow] (validated.south) to[bend right=10] node[right,font=\scriptsize]{expiry} (decommissioned.north);
\draw[arrow] (decommissioned) -- node[above,font=\scriptsize]{retention} (archived);
\end{tikzpicture}
\caption{TAC lifecycle state machine. Seven states with their
principal transitions. Each transition is a recorded
governance decision; the audit trail is preserved across
\emph{Decommissioned} and \emph{Archived}.}
\label{fig:cf-lifecycle}
\end{figure}

\paragraph{Governance over transitions.} Each transition is the
consequence of a recorded decision, not of an automated default.
The institutional registry that holds TAC declarations records
who validated, who activated, who suspended, who decommissioned;
the record is itself an audit artefact governed under the
canonical TAC schema of \trchap{ch:observability}{}. A TAC's
lifecycle is consequently fully reconstructable from the audit
trail.

\paragraph{Implications for procurement.} The lifecycle requires
that every component of the architecture support the full set
of transitions: provisioning on activation, release on
deactivation, preservation of classification metadata on
decommission, and read-only access to archived audit trails.
These properties are part of the procurement clauses developed in
\trchap{ch:cross-cutting}{} (\S\ref{sec:cc-procurement}).

%-----------------------------------------------------------------------------
\section{Out of scope}
\label{sec:cf-out-of-scope}

This Technical Reference does not address several adjacent
domains that an institution implementing the architecture will
need to engage separately.

\begin{itemize}[leftmargin=1.5em,nosep]
  \item \textbf{Human-resources policy.} Decisions about which
    subjects hold which roles, how role changes are processed in
    the institutional system of record, and how role lifecycle
    interacts with employment lifecycle are addressed by the
    institution's HR governance, not by this Reference.
  \item \textbf{Pedagogical training.} The training of academic
    staff and students in the practical use of TAC-aware
    services is an institutional learning-and-development
    function. The Reference's training discussion
    (\trchap{ch:cross-cutting}{} \S\ref{sec:cc-training}) addresses
    operational-team skills, not end-user pedagogy.
  \item \textbf{Procurement law.} The Reference proposes
    procurement language (\trchap{ch:cross-cutting}{}
    \S\ref{sec:cc-procurement}) consistent with the architecture
    but does not engage the specifics of any jurisdiction's
    procurement law. Institutions verify the language with their
    own legal counsel.
  \item \textbf{Physical security.} The Reference assumes that
    the physical infrastructure on which the architecture
    operates is secured to a level appropriate for its
    classification. Physical-security framework selection and
    operation is out of scope.
  \item \textbf{Disaster recovery beyond reversibility.} The
    reversibility playbook of \trchap{ch:reversibility}{}
    addresses migration-wave reversal. Broader disaster-recovery
    planning (data-centre loss, regional outage) is engaged by
    the institution's existing DR framework and is not
    superseded by this Reference.
\end{itemize}

\noindent The boundaries above are deliberate. The Reference is
intended to be useful in any institution's existing governance
ecosystem, not to displace that ecosystem.
